\documentclass[11pt]{amsart}

\usepackage[T1]{fontenc}
\usepackage{lmodern}
\usepackage{microtype}
\usepackage{amsmath,amssymb,amsthm}
\usepackage{booktabs}
\usepackage[hidelinks]{hyperref}

\hypersetup{
  pdftitle={A Counterexample to the Class 69 Equidistribution Conjecture on Involutions}
}

\newtheorem{theorem}{Theorem}
\theoremstyle{remark}
\newtheorem{remark}[theorem]{Remark}

\newcommand{\Sn}{\mathfrak{S}_n}
\newcommand{\In}{\mathcal{I}_n}
\newcommand{\occ}{\operatorname{occ}}

\title[Class 69 on involutions]
{A Counterexample to the Class 69 Equidistribution Conjecture on Involutions}
\date{}

\subjclass[2020]{05A05, 05A15}
\keywords{mesh pattern, permutation pattern, involution, equidistribution,
counterexample}

\begin{document}

\begin{abstract}
Fang, Fu, Kitaev, Li, Su, and Sun conjectured that the four length-two
mesh patterns in Class~69 are equidistributed on involutions. We
refute their conjecture, as stated in arXiv:2606.14367v1, at the
minimal size $n=3$. The two reverse-complement pairs have involution
distribution polynomials $2+q+q^2$ and $1+2q+q^2$, respectively; in
particular, their avoidance counts are different.
\end{abstract}

\maketitle

\section{The counterexample}

Mesh patterns were introduced by Br\"and\'en and Claesson~\cite{BC11}.
Using the convention that a mesh cell is indexed by the coordinates of
its lower-left corner, let $p_r=(12,R_r)$, where the four Class~69
patterns are listed in the order of Fang et al.~\cite[Conjecture~1]{FFKLSS26}:
\[
\begin{aligned}
R_1&=\{(0,0),(1,1),(1,2),(2,1)\},\\
R_2&=\{(0,1),(1,0),(1,1),(2,2)\},\\
R_3&=\{(0,2),(1,0),(1,1),(2,1)\},\\
R_4&=\{(0,1),(1,1),(1,2),(2,0)\}.
\end{aligned}
\]
Class~69 is the class of that name in \cite[Remark~1.3]{FFKLSS26},
described there as the only remaining conjectural equidistribution class
for mesh patterns of length two.

For completeness, let $p=(12,R)$ and $\sigma\in\Sn$. Given
$i<j$ with $\sigma_i<\sigma_j$, set
\[
 x_0=0,\quad x_1=i,\quad x_2=j,\quad x_3=n+1
\]
and
\[
 y_0=0,\quad y_1=\sigma_i,\quad y_2=\sigma_j,\quad y_3=n+1.
\]
The pair $(i,j)$ is an occurrence of $p$ if, for every $(a,b)\in R$,
there is no $k$ satisfying
\[
 x_a<k<x_{a+1}
 \qquad\text{and}\qquad
 y_b<\sigma_k<y_{b+1}.
\]
Here an occurrence is indexed by the pair of positions $(i,j)$, whereas
Fang et al.\ index it by the pair of values $(\sigma_i,\sigma_j)$; since
$i\mapsto\sigma_i$ is a bijection, the two conventions describe the same
occurrences and give the same count.
Write $\occ_p(\sigma)$ for the number of occurrences. For
\[
 \In=\{\sigma\in\Sn:\sigma^2=\mathrm{id}\},
\]
define
\[
 F^{\mathcal I}_{n,p}(q)
 =\sum_{\sigma\in\In}q^{\occ_p(\sigma)}.
\]
Equidistribution on involutions means equality of these polynomials
for every $n$.

The conjecture we refute is stated in the concluding remarks
of~\cite{FFKLSS26} and reads:
\begin{quote}
The patterns in the set $\{p_1,p_2,p_3,p_4\}$ are equidistributed on
involutions. (The first two patterns, as well as the last two patterns,
are trivially equidistributed via the composition of reverse and
complement.)
\end{quote}
In the source the four patterns are displayed as grid diagrams rather
than named; substituting $p_1,p_2,p_3,p_4$ for those diagrams, in the
order in which they appear there, is the only alteration we have made to
the quoted sentence.

\begin{theorem}
Conjecture~1 of Fang et al.~\cite{FFKLSS26}, as stated in
arXiv:2606.14367v1, is false. At $n=3$,
\[
 F^{\mathcal I}_{3,p_1}(q)
 =F^{\mathcal I}_{3,p_2}(q)
 =2+q+q^2,
\]
whereas
\[
 F^{\mathcal I}_{3,p_3}(q)
 =F^{\mathcal I}_{3,p_4}(q)
 =1+2q+q^2.
\]
\end{theorem}

\begin{proof}
The involutions in $\mathfrak S_3$ are
\[
 \mathcal I_3=\{123,132,213,321\}.
\]
For an increasing pair, the unique unused plot point lies in exactly
one mesh cell, and the pair is an occurrence of $p_r$ precisely when
that cell is not in $R_r$.

For $123$, the pairs $(1,2)$, $(1,3)$, and $(2,3)$ have unused-point
cells $(2,2)$, $(1,1)$, and $(0,0)$, respectively. For $132$, the
pairs $(1,2)$ and $(1,3)$ have cells $(2,1)$ and $(1,2)$. For $213$,
the pairs $(1,3)$ and $(2,3)$ have cells $(1,0)$ and $(0,1)$.
Finally, $321$ has no increasing pair. Hence
\[
\begin{array}{c@{\qquad}cccc}
\toprule
 &123&132&213&321\\
\midrule
\occ_{p_1}&1&0&2&0\\
\occ_{p_2}&1&2&0&0\\
\occ_{p_3}&2&1&1&0\\
\occ_{p_4}&2&1&1&0\\
\bottomrule
\end{array}
\]
and summing the corresponding powers of $q$ gives the stated
polynomials. In particular, the avoidance counts are $2$ for $p_1$
and $p_2$, but $1$ for $p_3$ and $p_4$.

The separation is already visible in a coarser invariant. The seven
increasing pairs listed above occupy seven distinct cells, namely every
cell except $(0,2)$ and $(2,0)$. Hence $R_1$ and $R_2$, which contain
neither of these two cells, each block four of the seven pairs, while
$R_3$ contains $(0,2)$ and $R_4$ contains $(2,0)$, so each blocks only
three. The total occurrence counts
$\sum_{\sigma\in\mathcal I_3}\occ_{p_r}(\sigma)$, that is, the values
$\frac{d}{dq}F^{\mathcal I}_{3,p_r}(q)\big|_{q=1}$, are therefore
$3,3,4,4$.
\end{proof}

\begin{remark}
The counterexample is minimal. For $n=0$ and $n=1$, no length-two
pattern can occur. For $n=2$, the involution $12$ has one occurrence
of every $p_r$, while $21$ has none, so
$F^{\mathcal I}_{2,p_r}(q)=1+q$ for all $r$.
\end{remark}

\begin{remark}
Zhang and Zhao~\cite[Theorem~1.1]{ZZ26} later proved the corresponding
Class~69 equidistribution on the full symmetric groups $\mathfrak S_n$.
The theorem above shows that this equivalence is not preserved after
restriction to involutions.
\end{remark}

\begin{thebibliography}{9}

\bibitem{BC11}
P.~Br\"and\'en and A.~Claesson,
\emph{Mesh patterns and the expansion of permutation statistics as
sums of permutation patterns},
Electron. J. Combin. \textbf{18} (2011), no.~2, Paper P5, 14~pp.,
doi:10.37236/2001.

\bibitem{FFKLSS26}
Q.~Fang, S.~Fu, S.~Kitaev, H.~Li, X.~Su, and Z.~Sun,
\emph{On mesh patterns of short length: Equidistribution and
enumeration},
arXiv:2606.14367v1 [math.CO], 2026.

\bibitem{ZZ26}
Z.-R.~Zhang and H.~Zhao,
\emph{The last distribution-equivalence class of mesh patterns of
length 2},
arXiv:2607.04111v1 [math.CO], 2026.

\end{thebibliography}

\end{document}